% ===================================================================
%  तक्ता क्रमांक ११ — शहर आणि त्याच्या वास्तू. --- आग.
%
%  Traduction, faite en 2026, en marathi d'aujourd'hui, sur l'ido de
%  texto/io. Regles de la colonne : voir 10-takta-01.tex. Deux
%  parties, deux numerotations : les cles portent c1 et c2. Le
%  feuillet 77 porte aussi la ligne d'apparat « तिसरी मालिका ».
%
%  LA VILLE MARATHE A SES MOTS, ET CE SONT DES MOTS DE VILLE
%  MARATHE : पेठ le quartier — celui des पेठा de Pune —, गल्ली la
%  ruelle, चौक la place, दरवाजा la porte de fort, बुरूज la tour,
%  तुरुंग la prison, घंटाघर le beffroi, घुमट la coupole, नगराध्यक्ष
%  le maire, जिल्हाधिकारी le prefet, पुतळा la statue, ग्रंथालय la
%  bibliotheque, तिवई le chevalet.
%
%  ET LE POMPIER SE NOMME PAR SA MACHINE. Le marathi appelle la
%  pompe a incendie बंब — le meme mot que le chaudron a eau chaude
%  du tableau 6, parce que c'est le meme objet : un reservoir de
%  cuivre qu'on chauffe ou qu'on vide. La pompe a vapeur (67) est
%  donc वाफेचा बंब, mot pour mot.
%
%  LE THEATRE EST ENTIEREMENT MARATHE, ce qui n'a rien d'un hasard :
%  le Maharashtra a une tradition de scene ancienne et vivante.
%  नाट्यगृह la salle, प्रयोग la representation — le mot ordinaire,
%  celui qu'on lit sur une affiche —, नट et नट्या les acteurs et les
%  actrices, वेशभूषा le costume, वाद्यवृंद l'orchestre, बासरी la
%  flute, ढोल le tambour.
%
%  ET LES BADAUDS SONT बघे, en un mot : ceux qui regardent. L'ido
%  doit dire « la kuriozi », les curieux ; le marathi a un nom pour
%  la foule qui s'attroupe devant un incendie.
%
%  LE RENVOI (81) EST POSE DEUX FOIS PAR L'IMPRIMEUR — sur les
%  musiciens et sur le violon. La traduction l'ecrit deux fois aussi :
%  la source ne bouge pas.
%
%  QUATRE INVERSIONS PREVENUES — la porte du chateau, la coupole de
%  l'hopital, l'incendie du theatre, les musiciens de l'orchestre —
%  ET TROIS SEULEMENT QUI ONT PASSE. C'est le meilleur rapport de la
%  colonne depuis le tableau 5, et la raison est simple : ce tableau
%  enumere des batiments plutot qu'il ne les rattache les uns aux
%  autres, et la postposition ne mord que sur « X de Y ».
%
%  ET UNE TROISIEME EXEMPTION DE LANGUE, LA PLUS ETROITE DES TROIS.
%  kolonoj.py relevait « दरवाजा » comme du hindi, et il a raison
%  partout ailleurs : le marathi dit दार pour la porte d'une maison,
%  et la colonne l'ecrit ainsi depuis le tableau 5. Mais दरवाजा est
%  le mot de la GRANDE porte d'un fort, celle qu'on appelle
%  महादरवाजा au Maharashtra, et c'est exactement ce dont parle
%  l'alinea 3 — la porte du vieux chateau fortifie, flanquee de ses
%  deux tours. Exemptee par sa forme exacte AVEC SON RENVOI :
%  verifie en posant un दरवाजा ordinaire au tableau 2, que le
%  controle releve aussitot.
% ===================================================================

%%K t11-apar-1 apar
\VUsaut{2.0mm}
\VUcentre{13.2pt}{90}{तिसरी मालिका}
\VUsaut{3.0mm}
\VUfilet{16mm}
\VUsaut{5.6mm}
\VUtitre{15.0pt}{60}{तक्ता क्रमांक ११}
\VUsaut{1.4mm}
\VUpk{11.4pt}{40}{शहर आणि त्याच्या वास्तू. --- आग.}
\VUsaut{2.2mm}

%%K t11-tit-1 sub
\VUpk{10.2pt}{40}{उपनगर.}
\VUsaut{3.0mm}

%%K t11-c1-01-1 p
१. --- मी महासागरापासून १०० किलोमीटरवर वसलेल्या एका मोठ्या शहरात राहतो, आणि तरीही ते पहिल्या
दर्जाचे \VUgras{बंदर} \textsuperscript{(1)} आहे. त्याच्या कडेने वाहणारी नदी ४०० मीटर रुंद आहे
आणि तिथे फार सुंदर नांगरतळ तयार होते.

%%K t11-c1-02-1 p
२. --- \VUgras{घाटांवरचा} \textsuperscript{(2)} शहराचा सगळा भाग पूर्णपणे सागरी आणि औद्योगिक
दिसतो, आणि तिथे एक \VUgras{उपनगर} \textsuperscript{(3)} तयार होते, ज्यात केवळ नावाडी आणि
कामगारच राहतात. हे कामगार \VUgras{कारखान्यांत} \textsuperscript{(4)} आणि
\VUgras{गॅसच्या कारखान्यात} \textsuperscript{(5)} काम करतात, ज्याचे उंच \VUgras{धुराडे}
\textsuperscript{(6)} आणि \VUgras{गॅसटाकी} फार दुरून दिसतात.

%%K t11-c1-03-1 p
३. --- मध्ययुगात शहराला तटबंदी होती, आणि तिच्या तटांचे काही अवशेष अजूनही सापडतात — विशेषतः तो
\VUgras{दरवाजा} \textsuperscript{(8)}, जुन्या \VUgras{तटबंदीच्या किल्ल्याचा}
\textsuperscript{(7)}, आपल्या दोन \VUgras{बुरुजांनी} \textsuperscript{(9)} वेढलेला, जे आता
\VUgras{तुरुंग} म्हणून वापरले जातात.

%%K t11-c1-04-1 p
४. --- फार जुन्या दिसणाऱ्या या सगळ्या \VUgras{पेठेत} फक्त एकच वास्तू तुलनेने आधुनिक आहे : ते
\VUgras{जकात कार्यालय} \textsuperscript{(10)}. ते घाट आणि \VUgras{गल्ली}
\textsuperscript{(11)} यांच्यामध्ये आहे, जी जुन्या \VUgras{नगरभवनाकडे} \textsuperscript{(12)}
जाते. ती शेवटची वास्तू सहज ओळखू येते, कारण तिचे अवाढव्य \VUgras{घंटाघर} \textsuperscript{(13)}
जुन्या शहरावर उंच उभे आहे.

%%K t11-c1-05-1 p
५. --- रस्त्याच्या दुसऱ्या बाजूला जकात कार्यालयासारख्याच शैलीत बांधलेली एक इमारत
\VUgras{उभी आहे} : तो \VUgras{शेअरबाजार} \textsuperscript{(14)}. तिचा एक दर्शनी भाग एका लहानशा
चौकाकडे पाहतो, जिथे \VUgras{जिल्हाधिकाऱ्यांचे कार्यालय} \textsuperscript{(15)} आहे. खरे तर
आमचे शहर राजधानी नसले तरी जिल्ह्याचे महत्त्वाचे मुख्यालय आहे.

%%K t11-tit-2 sub
\VUsaut{5.0mm}
\VUpk{10.2pt}{40}{शहराचा सर्वात उंच भाग.}
\VUsaut{3.4mm}

%%K t11-c1-06-1 p
६. --- बऱ्यापैकी मोठी शिबंदी अतिशय आरोग्यदायी अशा विस्तीर्ण \VUgras{छावण्यांत}
\textsuperscript{(16)} राहते; त्या \VUgras{नगरउद्यानाच्या} \textsuperscript{(17)} जाळीपर्यंत
पसरल्या आहेत. या उद्यानाकडे जाणारा \VUgras{रुंद रस्ता} \textsuperscript{(19)}
\VUgras{बचत बँकेच्या} \textsuperscript{(18)} मागून जातो आणि \VUgras{कॅथेड्रलच्या}
\textsuperscript{(20)} चौकात उतरतो.

%%K t11-c1-07-1 p
७. --- या सगळ्या वास्तू एका लहानशा टेकडीवर अर्धवर्तुळाकार पसरल्या आहेत, जिथे आमचे विस्तीर्ण
\VUgras{माध्यमिक विद्यालयही} \textsuperscript{(21)} आहे.

%%K t11-c1-07-2 p
मोठ्या रस्त्याला मिळणाऱ्या थोड्या उतरत्या वाटेने खाली उतरले, की \VUgras{न्यायालयापाशी}
\textsuperscript{(22)} पोहोचतो, ज्याच्यासमोर एक आल्हाददायक \VUgras{सार्वजनिक बाग}
\textsuperscript{(26)} पसरली आहे. हा \VUgras{बागेचा चौक} फार मोठा नाही, पण तो शहराच्या
मध्यभागी हिरवाईचे आणि गारव्याचे बेट तयार करतो. म्हणूनच शहरवासीयांची तिथे फार वर्दळ असते —
त्यांना \VUgras{कॅफेच्या} \textsuperscript{(25)} छत्रीखाली बसायला आवडते, हिरवळी आणि
फुलझाडांच्या ताटव्यांमध्ये उभ्या केलेल्या \VUgras{मंडपात} \textsuperscript{(27)} गुरुवारी आणि
रविवारी वाजवणारे लष्करी वादक ऐकायला.

%%K t11-c1-08-1 p
८. --- त्या हिरवळींपैकी एकीवर एक ब्राँझचा \VUgras{पुतळा} \textsuperscript{(28)} उभा आहे —
आपल्या जन्मगावाची मोठी सेवा करणाऱ्या आमच्या एका सहनागरिकाच्या स्मरणार्थ उभारलेले स्मारक.

%%K t11-tit-3 sub
\VUsaut{4.0mm}
\VUpk{10.2pt}{40}{वाहतुकीची साधने.}
\VUsaut{3.4mm}

%%K t11-c1-09-1 p
९. --- आतापर्यंत आमचे शहर विजेच्या दिव्यांनी उजळलेले नाही, त्याच्याकडे फक्त
\VUgras{गॅसचे दिवे} \textsuperscript{(29)} आहेत. पण \VUgras{इथली वृत्तपत्रे} ही प्रकाशव्यवस्था
सुधारावी म्हणून मोहीम चालवतात, \VUgras{ज्याप्रमाणे} \VUgras{ट्रामच्या ओढण्याची व्यवस्था} आधीच
सुधारली गेली. \VUgras{घोड्यांनी} \VUgras{ओढल्या जाणाऱ्या} \VUgras{त्या} जुन्या गाड्यांची जागा
प्रत्यक्षात \VUgras{विजेच्या ट्रामनी} \textsuperscript{(31)} घेतली, ज्या फार \VUgras{स्वस्त}
दरात प्रवाशांना \VUgras{शहराच्या एका टोकापासून दुसऱ्या टोकापर्यंत} भराभर नेतात.

%%K t11-c1-10-1 p
१०. --- न्यायालयाजवळ \VUgras{बँक} \textsuperscript{(32)} आहे, जिथे व्यापारी हुंड्या वटवल्या
जातात. \VUgras{बँकेचे वसुली कारकून} \textsuperscript{(33)} घरोघरी जाऊन येणी वसूल करतात.

%%K t11-tit-4 sub
\VUsaut{4.0mm}
\VUpk{10.2pt}{40}{कला आणि शिक्षण.}
\VUsaut{3.4mm}

%%K t11-c1-11-1 p
११. --- जवळच उभी असलेली ती सुंदर इमारत म्हणजे \VUgras{संग्रहालय}. त्यात शहराचे
\VUgras{ग्रंथालय} \textsuperscript{(34)}, \VUgras{निसर्गविज्ञानाचे} सुंदर \VUgras{संग्रह}
\textsuperscript{(35)} (निसर्गसंग्रहालय) आणि एक रेखीव \VUgras{चित्रसंग्रहालय}
\textsuperscript{(36)} एकत्र आहे, जे एक भव्य कायम \VUgras{प्रदर्शन} बनते.

%%K t11-c1-11-2 p
ते लोकांसाठी आठवड्यातून दोनदा उघडते, पण परगावच्या लोकांना \VUgras{रखवालदाराला}
\textsuperscript{(37)} बक्षिशी देऊन कोणत्याही दिवशी पाहता येते. \VUgras{चित्रकार}
\textsuperscript{(38)} तिथे काम करायला येतात. आपल्या तिवईसमोर, हातात रंगफलक घेऊन, प्रसिद्ध
चित्रांच्या नकला करताना ते दिसतात.

%%K t11-c1-12-1 p
१२. --- संग्रहालयामागच्या उंचवट्यावर \VUgras{घुमट} \textsuperscript{(40)} दिसू लागतो —
\VUgras{रुग्णालयाचा} \textsuperscript{(39)} — आणि \VUgras{अध्यापक विद्यालयाच्या}
\textsuperscript{(41)} इमारती, जिथे आमच्या नगरशाळांचे शिक्षक तयार होत असत. नाट्यगृहाच्या चौकात
\VUgras{टपाल कार्यालय} \textsuperscript{(43)} थाटले आहे, एका \VUgras{खाजगी घरात}
\textsuperscript{(42)}.

%%K t11-tit-5 sub
\VUsaut{5.0mm}
\VUpk{11.4pt}{40}{आग.}
\VUsaut{3.0mm}

%%K t11-tit-6 sub
\VUpk{10.2pt}{40}{आगीची हाक.}
\VUsaut{3.4mm}

%%K t11-c2-01-1 p
१. --- शहरभर एक भीतिदायक बातमी पसरते : नाट्यगृहात नुकतीच आग लागली! मी \VUgras{चौकात}
\textsuperscript{(96)} पोहोचतो त्या क्षणी \VUgras{पदपथावर} \textsuperscript{(46)} आणि
\VUgras{सडकेवर} \textsuperscript{(47)} गर्दी आधीच जमली आहे. \VUgras{टपाल कर्मचारी}
\textsuperscript{(44)} आणि \VUgras{पत्रे वाटणाऱ्या} \textsuperscript{(45)} टपाल कार्यालयातून
बाहेर पडतात. एका \VUgras{ट्रामचालकाला} \textsuperscript{(48)} शहराच्या \VUgras{रुग्णवाहिकेला}
\textsuperscript{(50)} वाट करून देण्यासाठी आपली गाडी थांबवावी लागली, जी \VUgras{शल्यचिकित्सक}
\textsuperscript{(52)} आणि \VUgras{परिचारक} \textsuperscript{(51)} घेऊन येते — \VUgras{जखमीवर}
\textsuperscript{(53)} उपचार करायला, ज्याला \VUgras{स्ट्रेचरवरून} \textsuperscript{(54)}
\VUgras{वृत्तपत्रांच्या मंडपाजवळ} \textsuperscript{(55)} आणले आहे.

%%K t11-c2-02-1 p
२. --- कवायतीहून परतणारी पायदळाची एक तुकडी \VUgras{शहरातल्या घोडेस्वार पोलिसांबरोबर}
\textsuperscript{(61)} व्यवस्था राखायला थांबली. \VUgras{घोडेस्वार}, ज्यांचे घोडे मागच्या
पायांवर उठतात, ते \VUgras{शिपायांपेक्षा} \textsuperscript{(57)} किंवा \VUgras{पोलिसांपेक्षा}
\textsuperscript{(63)} अधिक चांगल्या रीतीने \VUgras{बघ्यांना} \textsuperscript{(56)} मागे
रेटतात. एरवीही पोलिस फार गुंतलेले आहेत. त्यांच्यातला एक जण \VUgras{पोलिसांना}
\textsuperscript{(64)} दुसऱ्या \VUgras{अपघातग्रस्ताला} \textsuperscript{(65)} कुठे न्यायचे ते
सांगतो — शिडीवरून पडलेला एक धाडसी अग्निशामक.

%%K t11-c2-03-1 p
३. --- \VUgras{अधिकारी} \textsuperscript{(66)} येणारा \VUgras{वाफेचा बंब}
\textsuperscript{(67)} कुठे लावायचा ते नेमके सांगतो. त्या जबरदस्त यंत्राची वाट पाहत त्याने
घाईने एक \VUgras{हातपंप} \textsuperscript{(68)} लावून घेतला, ज्याची लांब \VUgras{नळी}
\textsuperscript{(69)} आगीवर लोटासारखे पाणी फेकते. सहा अग्निशामक तो चालवतात, तर \VUgras{तोटी}
\textsuperscript{(70)} धरलेला त्यांचा सोबती \VUgras{फवारा} \textsuperscript{(71)} आगीच्या
मध्यभागी वळवतो.

%%K t11-c2-04-1 p
४. --- नाट्यगृहाच्या दुसऱ्या बाजूला मोठी \VUgras{शिडी} \textsuperscript{(72)} टेकवली आहे.
तिच्यामुळे अग्निशामकांना काळ्या \VUgras{धुराच्या} \textsuperscript{(75)} भोवऱ्यांतून उसळणाऱ्या
ज्वाळांजवळ जाता येते.

%%K t11-tit-7 sub
\VUsaut{5.0mm}
\VUpk{10.2pt}{40}{शौर्याची कृत्ये.}
\VUsaut{3.4mm}

%%K t11-c2-05-1 p
५. --- \VUgras{आग} \textsuperscript{(74)} लागली \VUgras{नाट्यगृहाच्या} \textsuperscript{(73)}
प्रवेशदालनात, जिथे त्या \VUgras{संगीतिकेची} तालीम चालली होती, जिचा पहिला \VUgras{प्रयोग}
दुसऱ्या दिवशी होणार होता. सुदैवाने प्रेक्षागृहात थोडेच \VUgras{प्रेक्षक}
\textsuperscript{(76)} होते. ती बिचारी माणसे \VUgras{सज्ज्यावर} \textsuperscript{(77)}
आश्रयाला गेली, जिथे धाडसी \VUgras{अग्निशामक} \textsuperscript{(86)} शिडी आणि दोरखंड घेऊन
त्यांना वाचवायला येतात.

%%K t11-c2-06-1 p
६. --- \VUgras{स्तंभमंडपाखाली} \textsuperscript{(78)}, जिथे \VUgras{जाहिरातफलक}
\textsuperscript{(79)} प्रयोगाची माहिती देतो, तिथे \VUgras{वादक} \textsuperscript{(81)} —
\VUgras{वाद्यवृंदाचे} \textsuperscript{(80)} — आपली वाद्ये घेऊन पळताना दिसतात :
\VUgras{व्हायोलिन} \textsuperscript{(81)}, \VUgras{बासरी} \textsuperscript{(82)},
\VUgras{व्हायोलनचेलो} \textsuperscript{(83)}, \VUgras{ट्रॉम्बोन} \textsuperscript{(84)},
\VUgras{ढोल} \textsuperscript{(85)}, वगैरे. \VUgras{सामानाचा} \textsuperscript{(87)} काही भाग
वाचवण्याचा प्रयत्न चालला आहे.

%%K t11-c2-07-1 p
७. --- \VUgras{नटांना} \textsuperscript{(89)} आणि \VUgras{नट्यांना} \textsuperscript{(88)},
\VUgras{गायकांना} आणि \VUgras{गायिकांना} आपल्या नाटकी \VUgras{वेशभूषेतच}
\textsuperscript{(90)} पळावे लागले. तो हे कसे घडले ते एका \VUgras{पत्रकाराला}
\textsuperscript{(92)} समजावून सांगतो. \VUgras{वार्ताहर} अधिकाऱ्यांबरोबर पहिल्या क्षणापासूनच
धावत आला : \VUgras{जिल्हाधिकारी} \textsuperscript{(91)}, \VUgras{नगराध्यक्ष}
\textsuperscript{(93)}, \VUgras{पोलिस आयुक्त} \textsuperscript{(94)} आणि \VUgras{न्यायाधिकारी}
\textsuperscript{(95)}, जे जबाबदारीचा निवाडा करण्यासाठी चौकशी करतील.

\VUsaut{6.0mm}
\VUfilet{22mm}
